% F-M -- HPB Markup: the canonical example A (a firmware line) as text and
% as the diagram it produces. Illustrative data. Requires hpb.sty v1.1 and
% the listings language "hpb" (defined in main.tex).
% Column positions are those computed by @hpb/core 1.0 for this source
% (x positions divided by the column width): one column per event in
% statement order; the merge terminals of the branches closed by one
% "release ... merges" statement share one column just before the baseline;
% a dated baseline after a shared terminal column gets a small spacer. The
% document is pure history, so the renderer draws no NOW marker.
%   1 v1.0 | 2 I-21 | 3 X-101 | 4 terminal X-101 | 5 v1.0.1 | 6 X-102 |
%   7 T-8 | 8 terminals I-21, X-102 | 9.4 v1.1 | 10.4 I-22 | 11.4 X-103 |
%   12.4 terminals T-8, I-22, X-103 | 14.4 v1.2
\begin{figure}[tbp]
\centering
\begin{minipage}{0.66\linewidth}
\begin{lstlisting}[style=hpb,basicstyle=\ttfamily\footnotesize,numbers=left,numberstyle=\tiny\color{black!45},numbersep=6pt,xleftmargin=16pt]
hpb "Product A: firmware line"
line A "Product A"
release v1.0 @2026-01
improvement I-21 "Faster pairing"
issue X-101
release v1.0.1 merges X-101
issue X-102 "Clock drift"
task T-8 "Migrate logs"
release v1.1 @2026-03 merges I-21, X-102
improvement I-22 "Low-power mode"
issue X-103 "Menu translation"
release v1.2 @2026-05 merges T-8, I-22, X-103
\end{lstlisting}
\end{minipage}

\medskip
\hpbfit{%
\begin{hpbgraph}
  % lanes: improvements above, issues and tasks below the main line
  \hpblane{1}{imp}{Improvement (I)}{0}{15}
  \hpblane{-1}{iss}{Issue (X)}{0}{15}
  \hpblane{-2}{task}{Task (T)}{0}{15}
  % main line: pure history, no NOW marker
  \hpbmainline{0}{15}
  \hpborigin{0}{Product A}
  % released baselines
  \hpbrelease{1}{v1.0}
  \hpbrelease{5}{v1.0.1}
  \hpbrelease{9.4}{v1.1}
  \hpbrelease{14.4}{v1.2}
  % branches (fork col, lane, start col, end col, end, merge col)
  \hpbbranch[label={I-21 Faster pairing}]{imp}{1}{1}{2}{8}{merge}{9.4}
  \hpbbranch[label={X-101}]{iss}{1}{-1}{3}{4}{merge}{5}
  \hpbbranch[forkoff=0.6, label={X-102 Clock drift}]{iss}{5}{-1}{6}{8}{merge}{9.4}
  \hpbbranch[forkoff=0.3, mergeoff=0.3, label={T-8 Migrate logs}]{task}{5}{-2}{7}{12.4}{merge}{14.4}
  \hpbbranch[label={I-22 Low-power mode}]{imp}{9.4}{1}{10.4}{12.4}{merge}{14.4}
  \hpbbranch[mergeoff=0.6, label={X-103 Menu translation}]{iss}{9.4}{-1}{11.4}{12.4}{merge}{14.4}
  % ordinal time axis with the calendar ticks given by the @dates
  \hpbaxis{0}{15}{-2.75}
  \hpbtick{1}{Jan 2026}{-2.75}
  \hpbtick{9.4}{Mar}{-2.75}
  \hpbtick{14.4}{May}{-2.75}
\end{hpbgraph}}
\caption{HPB Markup in practice (illustrative). Top: twelve lines of text
that describe a firmware line. Bottom: the diagram they produce, drawn here
with \texttt{hpb.sty} at the column positions the renderer computes for
this source. Nobody types a coordinate: statement order gives the columns,
the branch type gives the lane, and each \texttt{release \dots{} merges}
closes the listed branches (\cref{rule:origin,rule:closure}). The terminals
closed by one release share a column; the history has no NOW marker.}
\label{fig:markup-example}
\end{figure}
